\section{计算图与反向传播}

这是本节课最核心的内容。反向传播是现代深度学习的\textbf{基石}——所有神经网络的训练都依赖于它。

\subsection{为什么需要反向传播}

要训练神经网络，我们需要计算损失函数 $L$ 对所有权重 $W_1, W_2, \ldots$ 的偏导数。直接对整个网络写出解析导数公式面临三个问题：

\begin{enumerate}
    \item \textbf{繁琐}：涉及大量矩阵运算的链式求导
    \item \textbf{不灵活}：改变损失函数或网络结构后需要重新推导
    \item \textbf{不可行}：对于上百层的网络，手动推导不现实
\end{enumerate}

\begin{figure}[H]
\centering
\includegraphics[width=0.95\textwidth]{slides-images/slide-048.jpg}
\caption{手动求导的问题：繁琐、不灵活、对复杂函数不可行\protect\footnotemark}
\end{figure}
\footnotetext{Slides 第49页。}

\subsection{计算图}

解决方案是将复杂函数分解为一系列\textbf{简单操作}（加法、乘法、指数等）的组合，用\textbf{计算图}表示。

\begin{figure}[H]
\centering
\includegraphics[width=0.95\textwidth]{slides-images/slide-049.jpg}
\caption{计算图：将复杂函数分解为简单操作的有向图\protect\footnotemark}
\end{figure}
\footnotetext{Slides 第50页。}

对于任意神经网络，计算图从输入数据和参数出发，经过一系列操作节点，最终得到损失值 $L$。

\begin{importantbox}{计算图的通用性}
计算图可以表示任意可微函数。无论网络多复杂——从简单的两层 MLP 到 Neural Turing Machine——都可以画成计算图。这使得反向传播成为一种\textbf{完全通用}的梯度计算方法。
\end{importantbox}

\subsection{反向传播的核心思想}

\textbf{前向传播}：从输入到输出，逐步计算每个节点的值。

\textbf{反向传播}：从输出到输入，逐步计算损失对每个节点的梯度。核心工具是\textbf{链式法则}。

\subsubsection{简单示例：$f(x, y, z) = (x + y) \cdot z$}

\begin{figure}[H]
\centering
\includegraphics[width=0.95\textwidth]{slides-images/slide-051.jpg}
\caption{反向传播示例：$f = (x+y) \cdot z$，从右向左逐步计算梯度\protect\footnotemark}
\end{figure}
\footnotetext{Slides 第52页。}

设 $x = -2, y = 5, z = -4$。

\textbf{前向传播}：
\begin{enumerate}
    \item $q = x + y = -2 + 5 = 3$
    \item $f = q \cdot z = 3 \times (-4) = -12$
\end{enumerate}

\textbf{反向传播}（从后向前）：

\begin{enumerate}
    \item $\frac{\partial f}{\partial f} = 1$（起始点）
    \item $\frac{\partial f}{\partial z} = q = 3$（乘法节点：对 $z$ 的局部梯度是 $q$）
    \item $\frac{\partial f}{\partial q} = z = -4$（乘法节点：对 $q$ 的局部梯度是 $z$）
    \item $\frac{\partial f}{\partial y} = \frac{\partial f}{\partial q} \cdot \frac{\partial q}{\partial y} = (-4) \times 1 = -4$（链式法则）
    \item $\frac{\partial f}{\partial x} = \frac{\partial f}{\partial q} \cdot \frac{\partial q}{\partial x} = (-4) \times 1 = -4$（链式法则）
\end{enumerate}

\begin{importantbox}{反向传播的两个关键概念}
\begin{itemize}
    \item \textbf{上游梯度}（Upstream Gradient）：从网络末端传回到当前节点的梯度，即 $\frac{\partial L}{\partial \text{output}}$
    \item \textbf{局部梯度}（Local Gradient）：当前节点输出对输入的导数，即 $\frac{\partial \text{output}}{\partial \text{input}}$
\end{itemize}
\textbf{反向传播规则}：\textbf{下游梯度 = 上游梯度 $\times$ 局部梯度}。每个节点只需要知道自己的局部梯度和收到的上游梯度，就可以计算出传给前一层的下游梯度。
\end{importantbox}

\subsection{Sigmoid 函数的反向传播}

考虑一个更复杂的例子：Sigmoid 函数 $\sigma(x) = \frac{1}{1 + e^{-x}}$。

\begin{figure}[H]
\centering
\includegraphics[width=0.95\textwidth]{slides-images/slide-054.jpg}
\caption{Sigmoid 计算图：分解为乘法、加法、取反、指数、取倒数五个基本操作\protect\footnotemark}
\end{figure}
\footnotetext{Slides 第55页。}

将 $\sigma(x)$ 分解为基本操作的链：
$$w_0 x_0 + w_1 x_1 + w_2 \xrightarrow{\text{取反}} -(\cdot) \xrightarrow{\exp} e^{(\cdot)} \xrightarrow{+1} (\cdot)+1 \xrightarrow{1/x} \frac{1}{(\cdot)}$$

每一步的局部梯度都很简单：
\begin{itemize}
    \item $1/x$ 的导数：$-1/x^2$
    \item $+c$ 的导数：$1$
    \item $e^x$ 的导数：$e^x$
    \item $-x$ 的导数：$-1$
    \item $cx$ 的导数：$c$
\end{itemize}

反向传播从末端的 $\frac{\partial L}{\partial L} = 1$ 开始，逐节点乘以局部梯度，最终得到 $\frac{\partial L}{\partial w_0}, \frac{\partial L}{\partial w_1}, \frac{\partial L}{\partial x_0}, \frac{\partial L}{\partial x_1}$。

\begin{figure}[H]
\centering
\includegraphics[width=0.95\textwidth]{slides-images/slide-055.jpg}
\caption{Sigmoid 反向传播的完整计算过程：从右向左逐步传递梯度\protect\footnotemark}
\end{figure}
\footnotetext{Slides 第56页。}

\begin{knowledgebox}{Sigmoid 导数的优雅形式}
Sigmoid 函数有一个美妙的性质：$\sigma'(x) = \sigma(x)(1 - \sigma(x))$。这意味着在实现中，如果前向传播已经计算了 $\sigma(x)$ 的值，反向传播时不需要重新计算，直接用前向值即可。这也是为什么在很多实现中会将 Sigmoid 作为一个\textbf{整体节点}而非拆分为基本操作。
\end{knowledgebox}

\subsection{常见操作的梯度模式}

\begin{figure}[H]
\centering
\includegraphics[width=0.95\textwidth]{slides-images/slide-056.jpg}
\caption{基本操作的梯度模式：加法分发梯度，乘法交换梯度，max 路由梯度\protect\footnotemark}
\end{figure}
\footnotetext{Slides 第57页。}

\begin{importantbox}{三种基本操作的梯度行为}
\begin{itemize}
    \item \textbf{加法门}（Add Gate）：\textbf{分发器}——上游梯度原样传递给两个输入。$\frac{\partial(x+y)}{\partial x} = 1$，$\frac{\partial(x+y)}{\partial y} = 1$
    \item \textbf{乘法门}（Multiply Gate）：\textbf{交换器}——每个输入收到的梯度是上游梯度乘以\textbf{另一个}输入的值。$\frac{\partial(xy)}{\partial x} = y$，$\frac{\partial(xy)}{\partial y} = x$
    \item \textbf{Max 门}（Max Gate）：\textbf{路由器}——梯度全部传给值较大的输入，另一个输入的梯度为零
\end{itemize}
\end{importantbox}

\subsection{多分支节点的梯度}

当计算图中一个变量被多个后续节点使用时（即一个节点有多条出边），反向传播时需要将各分支的梯度\textbf{相加}。

\begin{figure}[H]
\centering
\includegraphics[width=0.95\textwidth]{slides-images/slide-060.jpg}
\caption{多分支梯度：当一个变量参与多个后续计算时，各分支的梯度需要相加\protect\footnotemark}
\end{figure}
\footnotetext{Slides 第61页。}

这源自多元链式法则：如果 $z$ 通过 $f$ 和 $g$ 影响 $L$，则：
$$\frac{\partial L}{\partial z} = \frac{\partial L}{\partial f} \cdot \frac{\partial f}{\partial z} + \frac{\partial L}{\partial g} \cdot \frac{\partial g}{\partial z}$$

\begin{warningbox}{梯度相加，不是取平均}
当一个变量有多个使用者时，它对最终损失的影响是所有路径影响的\textbf{总和}。这是多元链式法则的直接结果。在实现自动微分系统时，一个常见 bug 是漏掉某条分支的梯度贡献。
\end{warningbox}

\subsection{向量化的反向传播}

前面的例子都是标量运算。在实际神经网络中，操作的对象是\textbf{向量和矩阵}。

\begin{figure}[H]
\centering
\includegraphics[width=0.95\textwidth]{slides-images/slide-062.jpg}
\caption{向量化反向传播：输入和输出都是向量时，局部梯度变成 Jacobian 矩阵\protect\footnotemark}
\end{figure}
\footnotetext{Slides 第63页。}

对于向量函数 $\mathbf{y} = f(\mathbf{x})$，局部梯度变成 \textbf{Jacobian 矩阵}：

$$J = \frac{\partial \mathbf{y}}{\partial \mathbf{x}} = \begin{bmatrix} \frac{\partial y_1}{\partial x_1} & \cdots & \frac{\partial y_1}{\partial x_n} \\ \vdots & \ddots & \vdots \\ \frac{\partial y_m}{\partial x_1} & \cdots & \frac{\partial y_m}{\partial x_n} \end{bmatrix}$$

反向传播规则变为：$\frac{\partial L}{\partial \mathbf{x}} = \frac{\partial L}{\partial \mathbf{y}} \cdot J$（上游梯度乘以 Jacobian）。

\begin{figure}[H]
\centering
\includegraphics[width=0.95\textwidth]{slides-images/slide-064.jpg}
\caption{ReLU 的向量化 Jacobian：对角矩阵，每个元素指示对应输入是否激活\protect\footnotemark}
\end{figure}
\footnotetext{Slides 第65页。}

\begin{knowledgebox}{逐元素操作的 Jacobian 是对角矩阵}
对于逐元素操作（如 ReLU、Sigmoid），输出的第 $i$ 个元素只依赖于输入的第 $i$ 个元素。因此 Jacobian 矩阵是\textbf{对角矩阵}——大量元素为零。在实现中，不需要显式构造完整的 Jacobian，只需对上游梯度做逐元素操作即可。
\end{knowledgebox}

\subsection{矩阵运算的反向传播}

\begin{figure}[H]
\centering
\includegraphics[width=0.95\textwidth]{slides-images/slide-067.jpg}
\caption{矩阵乘法 $Y = XW$ 的反向传播：$\frac{\partial L}{\partial X} = \frac{\partial L}{\partial Y} W^T$，$\frac{\partial L}{\partial W} = X^T \frac{\partial L}{\partial Y}$\protect\footnotemark}
\end{figure}
\footnotetext{Slides 第68页。}

对于全连接层的核心操作 $Y = XW$（$X \in \mathbb{R}^{N \times D}$，$W \in \mathbb{R}^{D \times M}$），反向传播公式为：

$$\frac{\partial L}{\partial X} = \frac{\partial L}{\partial Y} \cdot W^T \qquad \frac{\partial L}{\partial W} = X^T \cdot \frac{\partial L}{\partial Y}$$

\begin{importantbox}{矩阵反向传播的维度检查技巧}
利用\textbf{维度匹配}来验证公式的正确性：
\begin{itemize}
    \item $\frac{\partial L}{\partial X}$ 的形状必须与 $X$ 相同：$(N \times D)$
    \item $\frac{\partial L}{\partial Y}$ 的形状为 $(N \times M)$，$W^T$ 的形状为 $(M \times D)$
    \item $(N \times M) \cdot (M \times D) = (N \times D)$ \checkmark
\end{itemize}
同理可验证 $\frac{\partial L}{\partial W}$ 的维度正确性。这个技巧在调试反向传播实现时非常有用。
\end{importantbox}

\subsection{本章小结}

反向传播通过链式法则在计算图上从后向前逐节点传播梯度。每个节点只需知道自己的局部梯度和收到的上游梯度。关键模式：加法分发、乘法交换、max 路由、多分支相加。向量化时局部梯度变为 Jacobian 矩阵。

%% ========== 自动微分 ==========

